%!TEX root = ../_fyp.tex
\chapter{Recommendations}
\label{c-recommendation}

Based on the implementation experience and benchmark outcomes in Chapters~\ref{c-results} and~\ref{chap:discussion}, we offer recommendations for practitioners, system maintainers, and future researchers.

\section{Recommendations for Practitioners}

Students and data-science practitioners using Kaggle notebooks should select the interaction mode that matches task risk and intent:

\begin{itemize}
  \item Use \textbf{Ask mode} to understand unfamiliar notebooks, explain individual cells, or diagnose errors before editing. Benchmarks showed zero side effects and strong cell-level accuracy \parencite{vaithilingam2022expectation}.
  \item Use \textbf{Code mode} when ready to copy runnable cells into the editor. Verify placement suggestions against notebook flow before inserting \parencite{mcnutt2023notebookassistants}.
  \item Use \textbf{Agentic mode} for large batched edits when willing to review dispatched operations. The two-call protocol plans writes and runs together; users should inspect the batch before relying on outcomes \parencite{mozannar2024reading}.
\end{itemize}

For empty cells, prefer Ask or Code mode clarifying questions over immediate code generation. In evaluation, both modes deferred full scripts until intent was clear.

\section{Recommendations for System Improvement}

\begin{enumerate}
  \item \textbf{Implement ReAct-style feedback.} Add optional replanning after \texttt{run\_cell} using scrape-based execution observers, closing the gap between dispatch-only Agentic mode and full autonomous repair \parencite{yao2023react,shinn2023reflexion}.
  \item \textbf{Expand context packing.} Adopt hierarchical or retrieval-augmented cell selection for long notebooks to improve workflow-level Ask answers and Code placement \parencite{lewis2020rag,jiang2023structgpt}.
  \item \textbf{Harden DOM selectors.} Maintain scrape utilities against Kaggle front-end changes, following lessons from web data extraction research \parencite{chasins2018roussillon}.
  \item \textbf{Add human evaluation.} Supplement automated rubrics with expert grading for explanation quality and code correctness \parencite{chen2021codex,wang2021welldocumented}.
  \item \textbf{Live browser benchmarks.} Replicate harness scenarios on actual \texttt{/edit} sessions to measure end-to-end latency and user-perceived trust \parencite{mozannar2024reading}.
\end{enumerate}

\section{Recommendations for Future Research}

\begin{itemize}
  \item Compare GLM~4.7 with additional code models \parencite{openai2023gpt4,brown2020language} on identical frozen snapshots.
  \item Study learning outcomes when \acf{BSDS} students use Ask versus Code versus Agentic assistance during coursework \parencite{finnieansley2022codex}.
  \item Explore narrower tool APIs inspired by API-bank and Gorilla benchmarks \parencite{li2023apibank,patil2023gorilla} tailored to competition notebooks \parencite{quaranta2022bestpractices}.
  \item Investigate session-level memory across turns to reduce repeated query rounds in Agentic mode \parencite{park2023generative}.
\end{itemize}

These recommendations preserve the current system's strengths: grounded context, mode separation, and bounded Agentic dispatch. They also identify clear paths toward fuller autonomous notebook assistance in future iterations.

\section{Deployment and Maintenance Checklist}
\label{sec:deploy-checklist}

Teams deploying or extending the artefact should complete the following steps:

\begin{enumerate}
  \item Install the Chrome extension in developer mode and register the native host manifest (\texttt{register.ps1}).
  \item Configure \texttt{CEREBRAS\_API\_KEY} and \texttt{CEREBRAS\_MODEL} in \texttt{testing/host/.env}.
  \item Verify native messaging connectivity by opening a Kaggle \texttt{/edit} tab and sending an Ask-mode query.
  \item Enable Agentic mode only after confirming snapshot ingestion (\texttt{data/notebooks/live/} updates on scrape).
  \item Run smoke scripts under \texttt{testing/host/scripts/} after any \ac{DOM} selector change.
  \item Re-execute benchmark suites before thesis or demo updates (Appendix~\ref{app:reproduce}).
\end{enumerate}

\section{Risk Mitigation for Classroom Use}
\label{sec:classroom-risks}

Instructors adopting the assistant in \acf{BSDS} labs should communicate mode boundaries clearly:

\begin{itemize}
  \item \textbf{Assessment policy.} Code and Agentic modes can generate assignment solutions; courses should specify when AI assistance is permitted \parencite{finnieansley2022codex}.
  \item \textbf{Verification habit.} Students must run and inspect generated cells before submission; Ask mode supports understanding without automation.
  \item \textbf{Data privacy.} Notebook snapshots persist locally on the host machine; shared lab computers require session cleanup procedures.
  \item \textbf{API cost.} Agentic batches consume two \ac{API} calls per message; instructors may cap usage during introductory labs.
\end{itemize}

\section{Technology Roadmap}
\label{sec:roadmap}

A pragmatic three-stage roadmap extends the current prototype:

\begin{description}
  \item[Stage~1 (completed).] DOM scrape, three modes, two-call Agentic dispatch, harness benchmarks.
  \item[Stage~2 (near term).] Retrieval-augmented context packing; hardened selectors; live-browser benchmark replication.
  \item[Stage~3 (research).] Optional ReAct loop with scrape-based execution feedback; cross-model evaluation; learning-outcomes study.
\end{description}

Stage~2 items deliver the highest return for maintenance effort without altering the fundamental extension--host trust boundary established in this FYP.

\section{Summary Table of Recommendations}
\label{sec:rec-summary}

Table~\ref{tab:rec-summary} consolidates the recommendations in this chapter for quick reference during deployment or thesis defence.

\begin{table}[!ht]
\centering
\caption{Consolidated recommendations by audience.}
\label{tab:rec-summary}
\small
\begin{tabular}{|p{2.8cm}|p{5.2cm}|p{5.2cm}|}
\hline
\textbf{Audience} & \textbf{Priority action} & \textbf{Supporting evidence} \\ \hline
Students & Match mode to task risk & Table~\ref{tab:mode-guidance} \\ \hline
Instructors & Publish AI-use policy per mode & \ref{sec:classroom-risks} \\ \hline
Maintainers & Run benchmarks after DOM changes & Appendix~\ref{app:reproduce} \\ \hline
Researchers & Replicate on frozen snapshots & Appendix~\ref{app:snapshots} \\ \hline
Future devs & Stage~2: RAG context packing & \ref{sec:roadmap}, \ref{sec:future-react} \\ \hline
\end{tabular}
\end{table}
